\section{ELBO 分解与训练目标}

\subsection{从对数似然到 ELBO}

与 VAE 类似，扩散模型的训练目标是最大化数据的对数似然 $\log p_\theta(x_0)$。由于直接优化不可行，我们转而优化其下界——证据下界（ELBO）：

$$\log p_\theta(x_0) \geq \text{ELBO} = \mathbb{E}_{q(x_{1:T}|x_0)}\left[\log \frac{p_\theta(x_{0:T})}{q(x_{1:T}|x_0)}\right]$$

等价地，训练时我们最小化负 ELBO：

$$L = -\text{ELBO}$$

\begin{knowledgebox}{ELBO 与 VAE 的联系}
ELBO 的推导与 VAE 完全一致：通过引入变分后验 $q(x_{1:T}|x_0)$ 并利用 Jensen 不等式得到下界。不同的是，在扩散模型中，$q$ 是预定义的前向过程（不是学习得到的），因此 ELBO 的间隙（gap）由前向过程的设计决定。
\end{knowledgebox}

\subsection{ELBO 的逐项分解}

通过利用马尔可夫性质和贝叶斯规则，ELBO 可以分解为 $T+1$ 个项的和。具体地，负 ELBO 可写为：

$$L = \underbrace{D_{\text{KL}}\big(q(x_T|x_0) \,\|\, p(x_T)\big)}_{L_T\text{（先验匹配项）}} + \sum_{t=2}^{T} \underbrace{D_{\text{KL}}\big(q(x_{t-1}|x_t, x_0) \,\|\, p_\theta(x_{t-1}|x_t)\big)}_{L_{t-1}\text{（去噪匹配项）}} + \underbrace{\big(-\log p_\theta(x_0|x_1)\big)}_{L_0\text{（重建项）}}$$

各项的含义：
\begin{itemize}
    \item $L_T$（先验匹配项）：前向过程的终点分布 $q(x_T|x_0)$ 与先验 $p(x_T) = \mathcal{N}(0,I)$ 的 KL 散度。由于 $q$ 是预定义的且当 $T$ 足够大时 $q(x_T|x_0) \approx \mathcal{N}(0,I)$，此项近似为 0，\textbf{无需优化}。
    \item $L_{t-1}$（去噪匹配项）：这是核心项，要求学习的反向转移 $p_\theta(x_{t-1}|x_t)$ 匹配真实的后验转移 $q(x_{t-1}|x_t, x_0)$。
    \item $L_0$（重建项）：最后一步从 $x_1$ 重建 $x_0$ 的负对数似然。
\end{itemize}

\begin{importantbox}{训练的核心：去噪匹配项 $L_{t-1}$}
训练目标的关键是最小化去噪匹配项 $L_{t-1}$，即使反向转移分布 $p_\theta(x_{t-1}|x_t)$ 尽可能接近真实后验转移 $q(x_{t-1}|x_t, x_0)$。这要求我们首先推导出 $q(x_{t-1}|x_t, x_0)$ 的闭式表达。
\end{importantbox}

\subsection{真实后验转移分布的推导}

利用贝叶斯规则：

$$q(x_{t-1}|x_t, x_0) = \frac{q(x_t|x_{t-1}, x_0) \cdot q(x_{t-1}|x_0)}{q(x_t|x_0)}$$

由于三个分布 $q(x_t|x_{t-1})$、$q(x_{t-1}|x_0)$ 和 $q(x_t|x_0)$ 均为高斯分布，利用高斯分布的乘除运算规则，可以证明 $q(x_{t-1}|x_t, x_0)$ 也是高斯分布：

$$q(x_{t-1}|x_t, x_0) = \mathcal{N}\big(x_{t-1};\, \tilde{\mu}_t(x_t, x_0),\; \tilde{\beta}_t I\big)$$

其中方差和均值分别为：

$$\tilde{\beta}_t = \frac{(1 - \bar{\alpha}_{t-1}) \cdot \beta_t}{1 - \bar{\alpha}_t}$$

$$\tilde{\mu}_t(x_t, x_0) = \frac{\sqrt{\bar{\alpha}_{t-1}}\, \beta_t}{1 - \bar{\alpha}_t}\, x_0 + \frac{\sqrt{\alpha_t}\, (1 - \bar{\alpha}_{t-1})}{1 - \bar{\alpha}_t}\, x_t$$

\begin{warningbox}{后验依赖 $x_0$}
注意 $q(x_{t-1}|x_t, x_0)$ 的均值 $\tilde{\mu}_t$ 同时依赖于 $x_t$ 和 $x_0$。然而在反向过程（生成过程）中，我们\textbf{没有} $x_0$ 的信息——这正是需要神经网络来预测的原因。方差 $\tilde{\beta}_t$ 仅由噪声调度参数决定，不依赖 $x_0$，因此无需学习。
\end{warningbox}

\subsection{本章小结}

ELBO 分解将训练目标转化为在每个时间步最小化反向转移分布与真实后验转移分布之间的 KL 散度。真实后验转移 $q(x_{t-1}|x_t, x_0)$ 是高斯分布，其方差可以直接计算，均值是 $x_t$ 和 $x_0$ 的线性组合。

%% ========== 第四节：反向过程的参数化 ==========

